\documentclass[11pt,a4paper]{article}
\usepackage{fontspec}
\setmainfont{Noto Serif}
\setsansfont{Noto Sans}
\usepackage{babel}
\babelprovide[main,import]{russian}
\usepackage{amsmath,amssymb,mathtools,booktabs,geometry,fancyhdr,microtype,xcolor,hyperref}
\geometry{left=21mm,right=21mm,top=21mm,bottom=21mm}
\definecolor{ink}{HTML}{173D57}
\hypersetup{colorlinks=true,urlcolor=ink,linkcolor=ink,pdftitle={BGS-F. Ограничения на микроскопический вклад каналов Вейля}}
\pagestyle{fancy}\fancyhf{}\lhead{\small BGS-F}\rhead{\small Каналы Вейля на гейзенберговских временах}\cfoot{\thepage}
\setlength{\headheight}{15pt}\setlength{\parindent}{1em}\setlength{\parskip}{2.5pt}
\setlength{\emergencystretch}{3em}\tolerance=3000
\newcommand{\Tr}{\operatorname{Tr}}
\newcommand{\EE}{\mathbb E}
\newcommand{\R}{\mathbb R}
\newcommand{\Z}{\mathbb Z}
\newcommand{\C}{\mathbb C}
\newcommand{\si}{\mathrm{Sine}_2}
\newcommand{\eps}{\varepsilon}
\newcommand{\sinc}{\operatorname{sinc}}
\newcommand{\sect}[1]{\section*{\sffamily\color{ink}#1}}
\newcommand{\subt}[1]{\par\medskip\noindent\textbf{#1.}\quad}
\begin{document}
\noindent{\small\sffamily\color{ink} МАТЕМАТИЧЕСКАЯ ФИЗИКА\hfill 19 НОЯБРЯ 2026}
\begin{center}{\LARGE\sffamily Жёсткость карты и резонансы собственных фаз}\\[4pt]
{\large Точные результаты для нелинейной квантовой карты BGS-F}
\end{center}
\sect{Итог третьего этапа}
Пусть $U_N=\exp(i10^{-80}NA_N)M_N$ --- исходная карта. Сопряжение $X\mapsto U_NXU_N^*$ задаёт унитарную матрицу $W_N$ в ортонормированном базисе из $N^2$ операторов Вейля. Тогда $|\Tr U_N^m|^2=\Tr W_N^m$. Для веса $w$ исходного парного теста (P) доказано следующее утверждение: если $J_N$ --- произвольные множества индексов Вейля с $|J_N|=O(N)$, то по всем нечётным $N$
\[
 \boxed{\quad\frac1{N^2}\sum_{v\in J_N}
 \left|\sum_{m=1}^{2N-1} w(m/N)(W_N^m)_{vv}\right|\longrightarrow0.\quad} \tag{1}
\]
Допускается выбор $J_N$ с учётом спектра. Например, в эту оценку попадает любой один ряд индексов длины $N$, включая арифметически выделенный ряд. Внутри временной суммы сохраняются комплексные фазы. Оценка относится к диагональному разложению следа по каналам, а не к произвольному разбиению путей, проходящих через эти каналы.

Есть и равномерная поточечная оценка по времени. В пространстве бесследовых матриц сожмём $W_N$ на индексы $v=(a,b)$ с $0<|a|+\sqrt2|b|\le N/16$ (центрированные представители). Для результата $K_N$ имеем
\[
 \|K_N^{\lceil400\log N\rceil}\|_{2\to2}\le N^{-12}
 \qquad\text{при всех достаточно больших нечётных }N.                   \tag{2}
\]
Посредством детерминантного тождества (2) даёт абсолютную ошибку $O(N^{-8}\log N)$ для каждого полного $|\Tr U_N^m|^2$, $N\le m\le2N$, при исключении всех бесследовых замкнутых путей, проводящих более $\lceil400\log N\rceil$ последовательных шагов в указанной области. Константы здесь асимптотические; порог по $N$ не вычислен.

\subt{Результаты для всех степеней} Для $N\ne7$ и каждого $m\ge1$ никакая собственная непустая группа ненулевых вейлевых мод не переходит точно в группу той же размерности. Матрица вероятностей перехода строго сжимает ортогональное дополнение констант. Порядок любого точного резонанса собственных фаз, в том числе произведений их отношений, делит порядок $S$ по модулю $N$. Любая ненулевая щель короче $3/4$ средней имеет трансцендентное отношение собственных значений на концах. Доказательства даны в разделах 13--17. Случай $N=7$ вычислен точно.

\subt{Статус} Полный парный тест (P) и слабый предел $\si$ открыты. Разделы 1--7 доказывают (1),(2), разделы 8--9 исследуют арифметические возвраты. В разделах 10--12 выведены классическая жёсткость и степенное рассеяние внешних мод при фиксированной итерации. В конце сводного PDF приложены два предыдущих отчёта с собственной нумерацией.
\newpage
\sect{1. Генератор в базисе Вейля}
Пусть $e_N(x)=e^{2\pi i x/N}$ и $\bar2,\bar8$ --- соответствующие обратные по модулю нечётного $N$. Напомним матрицы
\[
 (T(a,b)u)(x)=e_N(b(x+\bar2a))u(x+a),\quad
 M_{xy}=N^{-1/2}e_N[-\bar8(3x^2-2xy+3y^2)].
\]
Пусть $r(q)=(1-\cos14\pi q)(1+\cos2\pi q)$, $e^{i\alpha}=(-3+4i)/5$. Если $c_{(a,b)}$ --- коэффициент при $e^{2\pi i(bq+ap)}$ функции
\[
 f_0=\cos14\pi q-1/10+r(q)(1/10+\sin(2\pi q)/100+\cos(2\pi p-\alpha)/100),
\]
то $A=\sum c_dT(d)$, $U=e^{i\tau N A}M$, $\tau=10^{-80}$. Всего 32 метки $d$, $|d_1|\le1$, $|d_2|\le9$, $\sum|c_d|\le593/500$. Для $\omega(d,v)=d_1v_2-d_2v_1$ и $S=\left(\begin{smallmatrix}3&4\\2&3\end{smallmatrix}\right)$ имеем
\[
 T(d)T(v)=e_N(\bar2\omega(d,v))T(d+v),\qquad MT(v)M^*=T(Sv).
\]
Эти тождества проверяются непосредственно в координатах. Используем скалярное произведение $N^{-1}\Tr X^*Y$, в котором $T(v)$ ортонормированы. На целочисленной решётке введём матрицу $G$, складывая вклады
\[
 G_{v+d,v}=-2\tau N c_d\sin(2\pi\bar2\omega(d,v)/N).
\]
Она периодична при одновременном сдвиге индексов на $N\Z^2$, косоэрмитова и ограничена для каждого $N$. Пусть $E=e^G$, $L=E\mathcal S$, где $\mathcal S e_v=e_{Sv}$. Периодизация по выходному индексу даёт в точности $W_N$: это следует почленно из ряда экспоненты и формулы для $i\tau N[A,T(v)]$. Метки, лежащие над нулевым классом, изолированы от остальных при действии $E$. В конечной матрице нулевая метка даёт отдельный собственный вектор с собственным числом 1.

\subt{Оценка дальнего перехода} Пишем $n(v)=|a|+\sqrt2|b|$, $u=a+\sqrt2 b$, $s=a-\sqrt2 b$. Тогда $n=\max(|u|,|s|)$ и $n(d)<14$. Поскольку $|\sin(2\pi\bar2 j/N)|=|\sin(\pi j/N)|$, сумма модулей одного шага $G$ из $v$ не больше $80\tau n(v)$: используем $|\omega(d,v)|\le10n(v)$. Каждый шаг увеличивает $n$ не более чем на 14. Для $b_0=2$ это даёт
\[
 \sum_w |E_{wz}|e^{2n(w-z)}
 \le \sum_{j\ge0}\frac{(80\tau e^{28})^j}{j!}
               \prod_{l=0}^{j-1}(n(z)+14l)
 =(1-1120\tau e^{28})^{-n(z)/14}\le e^{\epsilon n(z)},                \tag{3}
\]
где $\epsilon=10^{-64}$. Использованы биномиальный ряд и $-\log(1-x)\le2x$ при $0\le x\le1/2$. Эта оценка использует конкретный кик и остаётся справедливой при неограниченно большом $\tau N$.
\newpage
\sect{2. Убывание на целочисленной решётке}
Пусть $g_0(v)=|s|-|u|$, $\lambda=3+2\sqrt2$. При $S$ координаты $u,s$ умножаются соответственно на $\lambda,\lambda^{-1}$. Имеем $|g_0(w)-g_0(z)|\le2n(w-z)$ и
\[
 g_0(Sv)-g_0(v)\le-(1-\lambda^{-1})n(v).
\]
Отсюда при $v\ne0$ по (3)
\[
 \sum_w |L_{wv}|e^{g_0(w)-g_0(v)}\le e^{-0.8 n(v)}.                   \tag{4}
\]
Такой выбор веса связан с анизотропными нормами для классических гиперболических отображений [1]. В (4) оценена экспонента точного квантового коммутатора с модульными фазами.

На ромбе $0<n(v)\le R=N/16$ различим два сжатия: $K_0$ от поднятой матрицы $L$ и $K$ от конечной $W_N$. Оба суть сжатия унитарных операторов. По (4) $K_0$ сжимает взвешенную норму $\sum_v |f_v|e^{g_0(v)}$ в $q=e^{-0.8}$ раз. Кроме того,
\[
 \|K-K_0\|\le C e^{-N}.                                                  \tag{5}
\]
В самом деле, в одном элементе складываются дополнительные выходы $w+Nj$ при $j\ne0$. Имеем $n(w+Nj-Sv)\ge Nn(j)-7R$, $n(Sv)<6R$. Из (3) их сумма по модулю не выше $\sum_{j\ne0}e^{-2Nn(j)+14R+6\epsilon R}$. Переход к операторной норме стоит не более $N^2$. Поскольку $14/16<1$, итоговая оценка имеет запас в показателе.

Прямая замена взвешенной нормы на обычную на всём ромбе стоила бы $e^{2R}$. Для доказательства (2) отделим малое ядро $P$ с $0<n(v)\le l=100\log N$ от оболочки $Q$, $l<n(v)\le R$. Здесь $P,Q$ обозначают ортогональные проекторы внутри ромба. Нормы последующих матриц без индекса --- обычные операторные нормы в $\ell^2$.

Назовём переход в оболочке допустимым, если $n(w-Sv)\le n(v)/10$. Когда $|u(v)|\ge |s(v)|$, следующий вектор остаётся в этом конусе и его норма возрастает более чем в пять раз. Действительно, $|u(w)|\ge(\lambda-0.1)n(v)$ и $|s(w)|\le(\lambda^{-1}+0.1)n(v)$. В другом конусе норма убывает более чем вдвое на каждом шаге, пока конечный вектор остаётся в нём. Перейти из убывающего конуса в возрастающий можно только один раз. Значит, цепочка допустимых переходов внутри оболочки имеет длину менее $d=\lceil4\log N\rceil+4$: на две монотонные части приходится не более $2\log_2 N+2$ переходов.

Пусть $H=QKQ$, $H_0$ оставляет в $QK_0Q$ лишь допустимые элементы. Тогда $H_0^d=0$. Неприемлемый элемент $K_0$ по (3) имеет модуль не выше $e^{-0.19l}$. Значит, $\Delta:=\|H-H_0\|\le N^2e^{-0.19 l}+Ce^{-N}=O(N^{-17})$. Из $\|H\|\le1$ телескопированием получаем $\|H^d\|\le d\Delta(1+\Delta)^{d-1}=O(dN^{-17})$.
\newpage
\sect{3. Возврат в малое ядро и время выхода}
Для $j\ge1$ из (4),(5) следует
\[
 \|PK^jP\|\le (2l+1)e^{2l}q^j+Cj e^{-N}.                               \tag{6}
\]
Действительно, на $P$ множитель для перехода между взвешенными и невзвешенными нормами с двух концов не больше $e^{2l}$. Для преобразования $\ell^2$ во входную $\ell^1$ достаточно квадратного корня мощности ядра, не превосходящего $2l+1$. Внутренние шаги $K_0^j$ могут покидать малое ядро. Ошибка периодизации ограничена телескопированием сжатий в обычной норме.

Зафиксируем $B_0=\lceil400\log N\rceil$. Раскроем первые $d$ шагов по первому посещению $P$. В операторной записи это
\[
 K^d=K^d P+\sum_{i=1}^d K^{d-i}P K H^{i-1}Q+H^dQ .
\]
Здесь при $i=1$ степень ноль действует тождественно; проекторы в концах обеспечивают правильное ограничение. Аналогично раскроем последние $d$ шагов по последнему посещению $P$, применив ту же формулу к $K^*$. В $K^{B_0}=K^d K^{B_0-2d}K^d$ члены с пропущенным крайним посещением вместе оцениваются сверху $2\|H^d\|+\|H^d\|^2$: это следует заменой каждого крайнего множителя на $K^d-H^dQ$ или на соответствующий левый вариант. В остальных $(d+1)^2$ членах стоит $P K^jP$ с $B_0-2d\le j\le B_0$ в окружении сжатий нормы не больше единицы.

При этих длинах $e^{2l}q^j=O(N^{200-0.8(400-8)})$. Из (6) и предыдущей оценки для $H^d$ следует (2), даже с запасом на степени логарифма. В частности,
\[
 \|K^m\|\le e^{-m/70}\quad (m\ge2B_0),\qquad
 \sum_{j\ge B_0}\|K^j\|=O(N^{-12}\log N).                              \tag{7}
\]
Для первого неравенства используем $\lfloor m/B_0\rfloor\ge m/(2B_0)$ и $B_0\le401\log N$ при больших $N$; для второго складываем блоки по $B_0$ последовательных степеней. Вся аргументация справедлива для нечётных составных $N$, в том числе кратных 7. Никакая гипотеза о микроскопическом спектре или смешивании $|U_{xy}|^2$ не использовалась.

\subt{Смысл оценки} Нулевая метка отвечает постоянному оператору и была удалена до сжатия. Теорема ограничивает норму \emph{всей} суммы переходов, не покидающих ромб. Пространство имеет размер порядка $N^2$. Норму $W_N^m$ на полном пространстве по (7) оценивать нельзя: дополнение ромба участвует в унитарном возврате. В следующем разделе оно учитывается без потери фазы в местах выхода и входа.
\newpage
\sect{4. Внешние посещения и полный форм-фактор}
На бесследовой части пишем $V=W_N|_{I^\perp}=\left(\begin{smallmatrix}K&B\\ C&D\end{smallmatrix}\right)$ по ромбу и его дополнению. Для $|z|<1$ точная формула Шура имеет вид
\[
 \det(I-zV)=\det(I-zK)\det(I-zS(z)),\qquad
 S(z)=D+zC(I-zK)^{-1}B.
\]
Из унитарности $V$, положив $x=z(I-zK)^{-1}By$ и применив $V$ к $(x,y)$, выводим
\[
 I-S(z)^*S(z)=(1-|z|^2)B^*(I-\bar zK^*)^{-1}(I-zK)^{-1}B .
\]
Поэтому $S(z)$ есть сжатие. Положим $S_T(z)=D+\sum_{j=0}^{T-1}z^{j+1}CK^jB$ при $T=B_0$. По (7) для $|z|\le1$ ошибка не больше $C N^{-12}\log N$.

На круге $|z|\le e^{-1/N}$ имеем $\|(I-zS(z))^{-1}\|=O(N)$. Разность нормированных в нуле логарифмов двух внешних определителей ограничена $O(N^{-9}\log N)$: факторизация разности даёт $\log\det(I+E_z)$ с $\|E_z\|=O(N^{-11}\log N)$, а размерность не больше $N^2$. Применяя неравенство Коши к коэффициенту при $z^m$ и используя $e^{m/N}\le e^2$, для $N\le m\le2N$ получаем
\[
 |p_{m,N}|^2
 =1-m[z^m]\log\det(I-zS_T(z))+O(N^{-8}\log N).                        \tag{8}
\]
Здесь $|\Tr K^m|\le N^2 e^{-m/70}$ по (7). Единица есть вклад тождественного оператора. Для полной суммы (P) с $1\le m<N$ можно сохранить член $\Tr K^m$; его взвешенный вклад после нормировки $N^{-2}$ есть $O(N^{-2})$. Для этого вплоть до $2T$ применяем (4),(5) в виде $|\Tr K^m|\le q^{m-1}\sum_{v\ne0}e^{-0.8n(v)}+O(N^2m e^{-N})$, а далее суммируем по (7).

Разложение $-\log\det(I-zS_T)=\sum_{r\ge1}\Tr((zS_T)^r)/r$ показывает, что в (8) перечислены ровно пути с хотя бы одним внешним посещением и длиной каждого внутреннего участка в вершинах не более $T$. От внешней вершины до следующей допускается один шаг $D$ либо переход $CK^jB$, проходящий через $j+1$ внутренних вершин. Коэффициент $m/r$ возвращает выбор начала у исходного замкнутого слова: подсчитаем пары «путь с отмеченным началом, одна из $r$ внешних позиций» двумя способами. Это верно и для повторяющихся слов. Ошибка в (8) относится к полной знаковой сумме исключённых путей.

\subt{Оставшийся детерминант} Коэффициенты $S_T$ вычисляются по точному однократному пропагатору и степеням внутреннего сжатия до $O(\log N)$. Но коэффициент логарифма при $z^m$ содержит произведения этих матриц суммарной длины порядка $N$. Унитарность ограничивает $S$ в круге; величина суммарного следа таких произведений на нужном масштабе ещё не вычислена.
\newpage
\sect{5. Ни одна итерация не коммутирует с ненулевым сдвигом}
Здесь классическая матрица в координатах $q,p$ есть $B_c=\left(\begin{smallmatrix}3&-4\\-2&3\end{smallmatrix}\right)$. Как выведено в предшествующем отчёте, $U_N$ квантует $F=\phi^\tau\circ B_c$, где $\phi^t$ --- поток векторного поля $Y=-J\nabla g/(2\pi)$, $J=\left(\begin{smallmatrix}0&1\\-1&0\end{smallmatrix}\right)$. Физический символ $g$ получается из $f_0$ заменой $r(q)$ на $r(q+1/2)$ только в члене с $\cos(2\pi p-\alpha)$. Смещение, соответствующее $T(a,b)$, равно $(-a,b)/N$ на торе. Различия скалярных фаз не меняют рассматриваемых сопряжений.

Построим сопряжение $H(x)=x+h(x)$ с $F H=H B_c$, контролируя две моды $h$. На периодических векторных полях оператор $\mathcal L h=h-B_c h\circ B_c^{-1}$ обратим в $C^0$. После разложения на собственные направления $B_c$ его обратный равен сумме устойчивого ряда $\sum_{j\ge0} B_c^j P_s f\circ B_c^{-j}$ и неустойчивого ряда $-\sum_{j\ge1} B_c^{-j}P_u f\circ B_c^j$. Обе суммы геометрически сходятся и $\|\mathcal L^{-1}\|<5$ (нормы проекторов меньше 2). Обозначим $d_\tau=\phi^\tau-\mathrm{id}$ на плоском поднятии. Из $\|Y\|<20$, $\|DY\|<10^3$ следуют $\|d_\tau\|\le20\tau$, $\mathrm{Lip}(d_\tau)\le e^{10^3\tau}-1$ и $\|d_\tau-\tau Y\|\le10^4\tau^2$. Границы на $Y,DY$ даёт конечный ряд Фурье из раздела 1.

Уравнение $\mathcal L h=d_\tau(\mathrm{id}+B_c h\circ B_c^{-1})$ решается сжатием в $C^0$ и даёт $\|h\|\le100\tau$. Поскольку $\|B_c\|<8$, подстановка обратно даёт
\[
 \|h-\tau\mathcal L^{-1}Y\|\le10^8\tau^2.                               \tag{9}
\]
Отображение $H$ сюръективно по степени и инъективно по экспансивности $B_c$: из $H(x)=H(y)$ близость $B_c^j x,B_c^j y$ следует для всех целых $j$. Радиус $1/100$ заведомо годится: малые центрированные разности тогда следуют $B_c$ без целочисленного скачка, а ограниченный в обе стороны линейный гиперболический вектор нулевой.

У частот $e_1=(1,0),e_2=(0,1)$ иных представителей их $B_c^{\mathsf T}$-орбит в поддержке $\widehat g$ нет. Действительно, вторая координата всех ненулевых итераций этих двух частот по модулю не меньше 3 (из формулы для степеней $B_c$ либо рекурсии с характеристическим многочленом $X^2-6X+1$). В поддержке же она не превосходит 1. Значит, на обеих частотах коэффициент $\mathcal L^{-1}Y$ равен $P_s\widehat Y$. Для левого устойчивого ковектора $\ell_s=(1/2,1/\sqrt2)$ его модули равны соответственно $|\widehat g(e_1)|/\sqrt2$ и $|\widehat g(e_2)|/2$. Квадраты модулей $\widehat g$ равны $101/40000$ и $1/40000$. По (9) обе моды $h$ ненулевые. Следовательно, период $h(x+a)=h(x)$ возможен только при $a=0\pmod{\Z^2}$.

Если сдвиг $t_a$ коммутирует с $F^j$, $j\ge1$, то $H^{-1}t_aH$ коммутирует с $B_c^j$ и изотопен тождественному. Его поднятие $x+v(x)$ удовлетворяет $v(B_c^j x)-B_c^jv(x)=k\in\Z^2$. За вычетом постоянного решения, оставшаяся функция равна нулю по ограниченности в обоих направлениях. Получаем сдвиг $t_b$. Из $H t_b=t_a H$ интегрирование периодической части по площади даёт $a=b$ на торе и затем $h(x+b)=h(x)$. Отсюда $a=0$, что и требовалось.
\newpage
\sect{6. Спектральная мера одного канала}
Для $v\ne0$ обозначим через $\mu_{v,N}$ спектральную вероятностную меру $W_N$ в векторе $e_v=T(v)$ с нормировкой раздела 1; её моменты равны $(W_N^j)_{vv}$. Утверждается, что все частичные слабые пределы при произвольном выборе $v=v_N\ne0$ абсолютно непрерывны. Их плотности относительно нормированного углового Лебега ограничены числом $(1+q)/(1-q)$, $q=e^{-0.8}$. При $n(v_N)\to\infty$ единственный предел есть равномерная мера.

\subt{Макроскопическое смещение} Выберем подпоследовательность, на которой физическое смещение $a_N=(-v_1,v_2)/N$ сходится на торе к $a\ne0$. Для фиксированного $j\ge1$ диагональный коэффициент есть нормированный след
\[
 V_N=T(v)^* U_N^j T(v) U_N^{-j},\qquad
 G_a=t_{-a}F^j t_a F^{-j}
\]
(второе выражение --- предельное классическое отображение). По разделу 5 $G_a$ не тождественно. Это вещественно-аналитическое отображение, следовательно, множество его неподвижных точек имеет нулевую площадь: в локальных координатах достаточно свойства нулевого множества нетривиальной аналитической функции.

Приведём аргумент, переводящий это свойство в оценку следа без предположения о невырожденности неподвижных точек. Усредним $T(k)^* V_N T(k)$ по $k$ в квадрате $\{0,\ldots,L-1\}^2$. Все следы равны; квадрат модуля нормированного следа ограничен квадратом нормированной нормы Гильберта--Шмидта среднего. При раскрытии последней по парам $k,k'$ возникают нормированные следы $V_N^*T(k-k')V_N T(k-k')^*$ (или сопряжённые после перестановки меток). Центральные фазы сокращаются. При фиксированном $L$ формула Егорова на фиксированных временах и следовая формула для гладких символов дают в пределе сверху
\[
 \int_{\mathbb T^2}\prod_{d=1}^2
 \left|\frac1L\sum_{r=0}^{L-1}e^{2\pi i r (G_a(x)-x)_d}\right|^2 dx .
\]
Знаки и перестановка координат в спаривании меток Вейля не меняют произведения. При использовании обратного отображения интеграл тот же по сохранению площади. Для ясности, здесь Егорова достаточно применить к конечному набору гладких характеров: перенос сдвигом точен, а для кика символический коммутатор, делённый на $i\hbar$, сходится к скобке Пуассона. Последнее следует из формулы произведения Вейля и $\sin y=y+O(y^3)$ после суммирования быстро убывающих коэффициентов; интегрирование по фиксированному времени даёт остаток $o(1)$ в норме. Линейная часть точна. Входящие сдвиги $a_N\to a$ сохраняют равномерные нормы производных, а нормированный след произведения гладких символов сходится к интегралу по выделению нулевой моды. См. также [2]. Интеграл стремится к нулю при $L\to\infty$ по доминированной сходимости. Значит, $(W_N^j)_{vv}\to0$.

\subt{Малое смещение} Теперь $n(v)=o(N)$, а $j$ всё ещё фиксировано. Для $R_0=\max(\sqrt N,n(v))$ и $R_i=8^iR_0$ все эти радиусы $o(N)$ при $i\le j$. Матрица перехода из шара $n\le R_i$ за пределы $n\le R_{i+1}$ имеет норму $o(1)$: по (3) центральное смещение не менее $2R_i$, а дополнительные поднятия оцениваются тем же экспоненциальным хвостом. В частности, для центрированного выхода все его поднятия за пределами шара также имеют норму больше $R_{i+1}$. Хвост суммы по поднятиям ограничивается непосредственно (3), переход к норме стоит не более $N^2$. Телескопирование $j$ унитарных шагов позволяет ограничить их исходным ромбом $N/16$ с ошибкой $o(1)$. В нём (4),(5) дают $|(W_N^j)_{vv}|\le e^{-0.8n(v)}q^{j-1}+o(1)$.
\newpage
\sect{7. Что следует на микроскопическом масштабе}
Если $n(v)\to\infty$, обе альтернативы раздела 6 дают нулевые нетривиальные моменты и тем самым равномерную предельную меру. Для ограниченных $v\ne0$ моменты каждого частичного предела ограничены по модулю $q^{|j|}$ при $j\ne0$, включая отрицательные степени по сопряжению. Ряд Фурье плотности абсолютно сходится и даёт заявленную равномерную границу. Разбор подпоследовательностей покрывает все метки и нечётные $N$. В частности, для дуг длины $d_N\to0$
\[
 \sup_{v\ne0}\ \sup_{|I|\le d_N} \mu_{v,N}(I)\to0.                       \tag{10}
\]
Иначе сходящаяся подпоследовательность центров и мер дала бы атом положительной массы по теореме Портманто на замкнутых окрестностях предельного центра.

В собственном базисе $U_N u_k=e^{i\vartheta_k}u_k$ мера имеет вид
\[
 \mu_{v,N}=\frac1N\sum_{j,k}|\langle u_j,T(v)u_k\rangle|^2
              \delta_{\vartheta_j-\vartheta_k}.
\]
Отсюда, в частности, $\sup_{v\ne0}N^{-1}\sum_j |\langle u_j,T(v)u_j\rangle|^2\to0$ в любом собственном базисе. Кратные фазы добавляют к массе атома в нуле неотрицательные слагаемые и не мешают выводу. Оценка включает частоты порядка $N$.

Рассмотрим на окружности ядро $D_N(\theta)=N^{-1}\sum_{m=1}^{2N-1}w(m/N)e^{im\theta}$. Здесь $w=|(\phi*\phi)/(2I^2)|^2$, $\phi(t)=e^{-1/(1-t^2)}\mathbf1_{|t|<1}$, $I=\int\phi$. Ограниченность $w$ и её полной вариации, после суммирования по частям, дают $|D_N(\theta)|\le C\min(1,1/(N|\sin(\theta/2)|))$. Аналогичная оценка с $C_s$ верна для $N^{-1}(1-e^{-s/N+i\theta})^{-1}$ при фиксированном $s>0$. Отделим окрестность радиуса $N^{-1/2}$ вокруг нуля. Внутри применяем (10), снаружи --- верхнюю границу $O(N^{-1/2})$. Тем самым
\[
 \sup_{v\ne0,\varphi}\left|\int D_N(\theta-\varphi)d\mu_{v,N}\right|=o(1),
 \quad
 \sup_{v\ne0,\varphi}\frac1N\left|[(I-e^{-s/N-i\varphi} W_N)^{-1}]_{vv}\right|
 =o(1).                                                                    \tag{11}
\]
Суммирование первой оценки при $\varphi=0$ по $O(N)$ каналам с дополнительным множителем $1/N$ доказывает (1). Нулевой канал, если включён, отдельно стоит $O(1/N)$. Для абелева теста тот же результат следует из второй оценки с $2s$ и удалением нулевого по времени члена.

\subt{Дефицит точности} Всего каналов $N^2$. Для полной суммы кандидату на предел нужен контроль среднего вклада одного канала в (11) на масштабе $1/N$. В частности, (10) не различает варианты, когда масса дуги длины $1/N$ равна разным константам, делённым на $N$. После исключения разреженных семейств всё ещё требуется совместная оценка остальных фаз. Утверждение (1) также не исключает спектральный всплеск при отдельном $m$, распределённый по множеству каналов.
\newpage
\sect{8. Строгий контроль точных арифметических возвратов}
Полезно отдельно проверить, не даёт ли арифметика исходной карты скалярной степени при временах порядка $N$. Для всех нечётных $N>7$ и всех целых $m\ge1$ при заданном $\tau$ имеем
\[
 0<|\Tr U_N^m|<N .                                                         \tag{12}
\]
Если $M_N^m$ скалярна, число $|\Tr U_N^m|^2$ при таком кике даже трансцендентно. Доказательство использует теорему Линдемана--Вейерштрасса в форме линейной независимости экспонент попарно различных алгебраических показателей над полем алгебраических чисел [3].

Положим $p_m(t)=\Tr((e^{itA}M)^m)$. Собственные числа $A$ и его спектральные проекторы алгебраичны (для различных собственных чисел применим формулу Лагранжа). Поэтому $p_m(t)$ и, на вещественной оси, $F_m(t)=|p_m(t)|^2/N^2$ суть конечные суммы $c_\nu e^{it\nu}$ с алгебраическими коэффициентами и вещественными алгебраическими частотами. После группировки ненулевая такая функция не обращается в нуль в ненулевой алгебраической точке. Имеем $|p_m(0)|^2=\#\ker(S^m-I\bmod N)>0$, так как чистое сопряжение переставляет вейлев базис. Это даёт нижнее строгое неравенство при $t=\tau N$. Если $M^m$ не скалярна, то $F_m(0)<1$ и тот же аргумент применяется к $F_m-1$.

В скалярном случае обозначим $D_m=\sum_{r=0}^{m-1}M^r A M^{-r}$. Раскрытие произведения киков до второго порядка и цикличность следа дают
\[
 F_m(t)=1-t^2\big(\Tr(D_m^2)/N-(\Tr D_m/N)^2\big)+O_{m,N}(t^3).
\]
Вариация равна $m^2\sum_{O\ne\{0\}}|\sum_{d\bmod N\in O} c_d|^2/|O|$, где $O$ --- орбита $S$ по меткам. Действительно, $S^m=I$ по ковариантности, и каждая орбита равномерно обходится. У метки $(0,7)$ ненулевая орбита; свой коэффициент равен $0.45$, коэффициент $(0,-7)$ тоже положителен, а сумма модулей остальных не больше $1.186-0.9=0.286$. Поэтому вариация положительна даже при наложении малых частот по модулю. Функция $F_m$ непостоянна; отняв произвольное алгебраическое число, применяем указанную линейную независимость. Получаем (12).

В частности, при $N=7^j$, $j\ge3$, $m=3N/7$ точная кривизна имеет вид
\[
 -F_m''(0)/(2m)=11370883/4000000.                                         \tag{13}
\]
Здесь $S^3-I=14\left(\begin{smallmatrix}7&10\\5&7\end{smallmatrix}\right)$. Для примитивной по 7 метки длина орбиты $m$, для двух делимых на 7 меток $(0,\pm7)$ она $m/7$. Для проверки расщепим $a+b\sqrt2$ по двум корням $r_\pm$ уравнения $r^2=2$ в $\Z_7$, равным $3,4$ по модулю 7 (простое поднятие Гензеля). В этих координатах $S$ умножает на взаимно обратные $3+2r_\pm$. Они имеют порядок 3 по модулю 7, а $(3+2r_\pm)^3-1=7(14+10r_\pm)$ имеет 7-адическую нормировку ровно один. Биномиальная формула даёт указанные длины и образ в каждой единичной координате --- полный прообраз квадратов $\bmod 7$.
\newpage
\sect{9. Орбитный расчёт и конечные проверки}
Для (13) достаточно посчитать орбиты 32 мод при модуле 343. Различные орбиты при 343 не сливаются при поднятии. Если две примитивные моды лежат в одной такой орбите, их целые нормы $a^2-2b^2$ равны: их разность делится на 343 и по абсолютной величине меньше 343. В координате, являющейся единицей, отношение сравнимо с квадратом по модулю 7 и потому принадлежит группе $3+2r_\pm$ на любой глубине. Другая координата определяется точно по норме и единичной координате. Две непримитивные метки лежат в разных орбитах, так как $-1$ не квадрат по модулю 7. Это доказывает стабилизацию разбиения на всех $j\ge3$.

Если $d=(a,b)$, то $(1,b),(-1,b)$ соединяются для $b=1,-1,6,-6,8,-8$; ещё соединяются $(a,7),(a,-7)$ для $a=1,-1$. Остальные классы одноэлементны на исходном носителе. В формуле вариации веса $m/|O|$ равны 1 и 7. Сложение квадратов гауссовых целых со знаменателем $4000^2$ даёт (13); при модуле 49 получается $11370483/4000000$. Формула Тейлора описывает только производную при нуле. При заданном $\tau N$ из этого расчёта выводится строгая нескалярность, но не равномерное расстояние от скалярной матрицы.

\subt{Независимые конечные проверки} Матрица $A$ собрана по трём диагоналям и отдельно по 32 модам. Сопряжение непосредственно вычислено на всех матрицах Вейля и сопоставлено с экспонентой генератора из раздела 1. Для проверки алгебраических тождеств использованы также увеличенные тестовые значения силы и отдельная тестовая маска, чтобы переходы между блоками были видны при малых $N$. Они не входят в доказательство при $10^{-80}$. В последнем столбце сравниваются коэффициент логарифма и прямой перебор циклических двоичных слов с ограниченной длиной внутреннего участка.
\begin{center}\small\begin{tabular}{@{}rrrr@{}}\toprule
$N$ & Ошибка сопряжения & Ошибка формулы Шура & Ошибка подсчёта слов \\ \midrule
5 & $4.0\cdot10^{-16}$ & $8.1\cdot10^{-16}$ & $2.3\cdot10^{-16}$ \\
7 & $5.3\cdot10^{-16}$ & $1.4\cdot10^{-17}$ & 0 \\
9 & $1.1\cdot10^{-15}$ & $5.6\cdot10^{-16}$ & $1.4\cdot10^{-15}$ \\ \bottomrule
\end{tabular}\end{center}
Константы для ухода, квадраты двух первых мод $g$ и орбитные суммы проверяются рациональной арифметикой. В архиве есть программа, исходные данные предыдущих этапов и инструкции воспроизведения. Аналитические рассуждения о пределе приведены в тексте.

\subt{Граница арифметической оценки} Остаточный объект теперь можно задать либо полным внешним детерминантом в (8), либо коллективной суммой спектральных мер оставшихся каналов в (11). Для первого объекта неизвестен асимптотический коэффициент при $m\asymp N$; для второго неизвестна средняя микроскопическая плотность с точностью $o(1/N)$ на канал после нужного сглаживания. В частности, даже строгая трансцендентность в (12) не даёт оценки величины этого остатка.

\subt{Литература} [1] F. Faure, N. Roy, \href{https://doi.org/10.1088/0951-7715/19/6/002}{Ruelle--Pollicott resonances for real analytic hyperbolic maps} (2006). [2] R. Schubert, \href{https://arxiv.org/abs/math-ph/0503045}{Upper bounds on the rate of quantum ergodicity} (2006), случай возмущённых карт тора. [3] A. Baker, \emph{Transcendental Number Theory}, Cambridge University Press, гл.~1, теорема Линдемана--Вейерштрасса. Обозначения BGS-F и (P) взяты из исходной постановки, приложенной к архиву.
\newpage
\sect{10. Жёсткость всех итераций и распространение мод}
Усилим раздел 5. Для каждого $j\ge1$ имеет место
\[
 F^j t_a F^{-j}=t_b\quad\Longrightarrow\quad a=b=0\quad\text{на торе}.  \tag{14}
\]
Здесь разрешены \emph{разные} сдвиги. В частности, классическая карта не имеет точного направления распространения трансляций даже после сколь угодно большого фиксированного числа итераций.

Положим $h_1=\mathcal L^{-1}Y$, $k=H^{-1}-\mathrm{id}$. Из (9)
\[
 \|k+\tau h_1\|_\infty\le10^8\tau^2+10^5\tau^{3/2}.
\]
Действительно, $k=-h\circ H^{-1}$, $\|k\|\le100\tau$, $[h_1]_{1/2}<10^4$: в двух рядах раздела 5 используем $[Y]_{1/2}\le10^3$, $\|B_c^{\pm n}\|\le2\lambda^n$ и суммируем по степеням $\lambda^{-1/2}$. Для $\ell_u=(1/2,-1/\sqrt2)$ и $k_i=B_c^{\mathsf T}e_i$ коэффициент $\ell_u\widehat h_1(k_i)$ равен $-\lambda^{-1}\ell_u\widehat Y(e_i)$, поскольку из всей орбиты в носитель попадает только $e_i$. Из чисел раздела 5 получаем
$|\widehat{\ell_u k}(k_i)|>10^{-4}\tau$ для обоих $i$.

Пусть $\Phi_a=H^{-1}t_aH$. При $x=H^{-1}(y)$ имеем для совместимых поднятий
$\ell_u(\Phi_a(x)-x)=\ell_u a+\ell_u k(y+a)-\ell_u k(y)$.
Если $a\ne0$, заменой обоих сдвигов одинаковой целой кратностью можно добиться $\|a\|_{\mathbb T,\infty}\ge1/4$: ненулевая подгруппа окружности в хотя бы одной координате встречает $[1/4,3/4]$. Тогда $\max_i\|k_i\cdot a\|_{\R/\Z}\ge1/28$, поскольку $\|B_c^{-1}\|_\infty=7$. Одна из двух выделенных мод даёт осцилляцию $\ell_u(\Phi_a(x)-x)$ не менее $10^{-5}\tau$. Здесь использованы $|e^{2\pi iu}-1|\ge4\|u\|_{\R/\Z}$ и то, что модуль ненулевого коэффициента не превосходит осцилляции. С другой стороны, для любого $b$ осцилляция не превосходит $400\tau$. При предполагаемом равенстве (14) имеем $B_c^j\Phi_a B_c^{-j}=\Phi_b$; неоднозначность поднятия даёт лишь константу. Осцилляция умножается на $\lambda^j$, откуда $\lambda^j\le4\cdot10^7$ и $j\le10$.

Для $1\le j\le10$ периодическая часть поднятия имеет разложение
\[
 F^j(x)-B_c^jx=\tau\sum_{i=1}^j B_c^{j-i}Y(B_c^ix)+E_j(x),
 \qquad \|E_j\|\le10^{20}\tau^2 .
\]
Действительно, $\|F^i-B_c^i\|\le20\tau(8^i-1)/7$ для поднятий, а на одном шаге линейная аппроксимация кика стоит $10^4\tau^2$ плюс $10^3\tau\cdot8\|F^{i-1}-B_c^{i-1}\|$. Суммируем с множителями $8^{j-i}$. В частоте $k_l$ выживает только $i=1$, с коэффициентом $B_c^{j-1}\widehat Y(e_l)$. Его норма больше $10^{-11}$, используя $\|B_c^{-1}\|<8$. По малости остатка обе моды ненулевые. Если $F^jt_a=t_bF^j$, интегрирование по площади сначала даёт $b=B_c^ja$ на торе, затем $a$ должен быть периодом этой периодической части. Моды $k_1,k_2$ составляют целочисленный базис. Значит, такой ненулевой $a$ невозможен. Это завершает доказательство (14).

\newpage
\sect{11. Степенная оценка для каждого внешнего перехода}
Для любых фиксированных $j\ge1$, $\delta>0$ существуют $C_{j,\delta}<\infty$ и $\nu_{j,\delta}>0$ такие, что
\[
 \boxed{\sup_{n(v)\ge\delta N}\ \sup_w |(W_N^j)_{wv}|
       \le C_{j,\delta}N^{-\nu_{j,\delta}}.}                              \tag{15}
\]
Представители центрированы. Это поточечное утверждение о распространении по каналам. Оно включает все выходные индексы и не требует никакой случайности фаз. Показатель может зависеть от $j$; на времена $j\asymp N$ (15) не переносится.

Вместо диагонального коэффициента раздела 6 берём нормированный след $T(w)^*U^jT(v)U^{-j}$. Его классическая карта $G=t_{-b}F^jt_a F^{-j}$. Параметры лежат в компакте, где $a$ отделён от нуля. По (14) ни один $G$ не тождествен. Существуют целое $D\ge1$ и $C$ такие, что равномерно на этом компакте
\[
 \operatorname{area}\{x:\operatorname{dist}(G(x),x)<\epsilon\}
 \le C\epsilon^{1/D}.
\]
Для проверки в окрестности каждой точки равенства $G(x)=x$ выбираем локальное поднятие одной компоненты разности и направление, в котором её производная некоторого конечного порядка отлична от нуля. Это возможно по аналитичности: иначе все ряды обнулились бы и карта была бы тождественна. На конечном покрытии по $x,a,b$ получаем нижние границы для этих производных, а вне покрытия разность отделена от нуля. На малом прямоугольнике интегрируем вдоль выбранного направления. Если множество $|f|<\epsilon$ на отрезке имеет длину $l$, в нём можно выбрать $d+1$ точек с расстояниями $\ge c_d l$. Теорема о разделённых разностях ограничивает значение $|f^{(d)}|$ в одной из промежуточных точек числом $C_d\epsilon/l^d$. Это даёт нужную оценку на каждом срезе.

В аргументе усреднения из раздела 6 допустим $L=\lfloor N^{1/24}\rfloor$. Ошибка для пары характеров не более $C_jN^{-1}(1+L)^{12}$, равномерно по сдвигам. Вот детали контроля порядка. На фиксированном числе гладких классических шагов $C^k$-норма перенесённого характера есть $O_{j,k}((1+L)^k)$. Ошибка коммутатора кика после деления на $\hbar$ ограничена $C N^{-2}\sum_l(1+|l|)^3|\widehat f(l)|\le C N^{-2}\|f\|_{C^7}$ по модульному правилу в физическом базисе. При перемножении квантований двух тестов ошибка в норме не превосходит $C N^{-1}\|f\|_{C^5}\|h\|_{C^5}$ по неравенству для фазы Вейля. При взятии нормированного следа ошибка от ненулевых мод в $N\mathbb Z^2$ ограничена $C N^{-4}\|fh\|_{C^4}$. Эти оценки, интегрирование по времени кика и конечное телескопирование дают заявленный грубый порядок 12. Веса при усреднении неотрицательны с суммой один.

Вне окрестности радиуса $L^{-1/2}$ хотя бы один нормированный геометрический ряд имеет квадрат модуля $O(1/L)$. Внутри применим оценку площади. Итак, квадрат левой части (15) не больше $C_{j,\delta}(L^{-1/(2D)}+L^{-1}+N^{-1}(1+L)^{12})$. Это доказывает (15), например с $\nu=1/(96D)$.

\newpage
\sect{12. Когерентные пакеты и нерешённая часть}
Пусть $S_N$ содержится во внешней области $n(v)\ge\delta N$ и $|S_N|=N^{o(1)}$. Для \emph{любого} единичного вектора $b_N$, поддержанного на $S_N$, спектральная мера $W_N$ в $b_N$ слабо сходится к равномерной мере окружности. Это включает произвольные сложные фазы коэффициентов самого пакета. Для фиксированного $j\ge1$ норма $P_{S_N}W_N^j P_{S_N}$ ограничена по (15) числом $C_{j,\delta}|S_N|N^{-\nu_{j,\delta}}\to0$. Отрицательные степени учитываются сопряжением. Теорема следует из плотности тригонометрических полиномов; сходимость равномерна по всем векторам в данной поддержке.

Отсюда следует утверждение об \emph{отдельных} собственных операторах. Для каждого $\varepsilon>0$ существует $\beta(\delta,\varepsilon)>0$ такое, что любое семейство $S_N\subset\{n\ge\delta N\}$ размера не более $N^\beta$ удовлетворяет при достаточно больших $N$
\[
 \sup_{W_N X=e^{i\theta}X,\,\|X\|_{\rm HS,n}=1}
              \sum_{v\in S_N}|\langle T(v),X\rangle_{\rm HS,n}|^2
 \le\varepsilon .
\]
Здесь $\|X\|_{\rm HS,n}^2=N^{-1}\Tr X^*X$, собственная фаза произвольна, включая кратные. В самом деле, для спектрального проектора $\Pi_\theta$ и единичного $b$ на $S_N$ имеем
\[
 \|\Pi_\theta b\|^2\le
 \left\|T^{-1}\sum_{j=0}^{T-1}e^{-ij\theta} W_N^jb\right\|^2
 \le T^{-1}+\frac{2}{T^2}\sum_{j=1}^{T-1}(T-j)
                       \|P_{S_N}W_N^jP_{S_N}\|.
\]
Берём $T>2/\varepsilon$, затем $0<\beta<\min_{1\le j<T}\nu_{j,\delta}$. Нормы $\|\Pi_\theta P_{S_N}\|$ и $\|P_{S_N}\Pi_\theta\|$ равны. Это даёт требуемое утверждение. Тем же методом, либо по (10) для меняющегося вектора, пакет из $N^{o(1)}$ мод не концентрируется спектрально в дугах длины $o(1)$.

\subt{Значение для BGS-F} Жёсткость (14) и дисперсия (15) запрещают перенос отдельной внешней моды или когерентного пакета малой размерности в одну из немногих выходных мод с ненулевой предельной амплитудой при фиксированной итерации. Самосогласованный вклад большого числа рассеянных мод при растущих временах этим не определяется. Константы (15) не дают суммируемой оценки для полного детерминанта (8) на гейзенберговской полосе. Именно этот предел, а также высшие факториальные порядки остаются открытыми.


\newpage
\sect{13. Связность при всех степенях точной матрицы}
Пусть $N\ge3$ нечётно, $N\ne7$, $m\ge1$ произвольно. На ненулевых метках образуем матрицу
\[
 Q_{wv}=|(W_N^m)_{wv}|^2,\qquad v,w\in(\Z/N\Z)^2\setminus\{0\}.
\]
Она дважды стохастична. Для \emph{каждой} такой пары $(N,m)$ докажем строгое неравенство для её второго сингулярного числа:
\[
                        \boxed{\sigma_2(Q)<1.}                            \tag{16}
\]
В частности, никакое собственное подмножество ненулевых мод не порождает инвариантное координатное подпространство для $W_N^m$. У (16) есть точное исключение: $A_7=(9/10)I$, поскольку $r(x/7)=0$ на всей решётке, включая сдвинутые модульные середины, а $\cos(14\pi x/7)=1$. Здесь $Q$ является перестановкой.

Обозначим $U(t)=e^{itA}M$, $P=S^m$, $D_m=\sum_{j=0}^{m-1}M^jAM^{-j}$. Для вещественного $t$ все коэффициенты матрицы $\operatorname{Ad}(U(t)^m)$ суть конечные экспоненциальные полиномы с алгебраическими частотами и коэффициентами. По теореме Линдемана--Вейерштрасса такой полином, ненулевой как функция, имеет ненулевое значение в каждой ненулевой вещественной алгебраической точке. Здесь используется линейная независимость экспонент \emph{различных} алгебраических чисел над $\overline{\mathbb Q}$. В частности, положителен переход $v\to x=Pv$, ненулевой при $t=0$. Производная образа $T(v)$ при нуле равна $i[D_m,T(x)]$. Если коэффициент $D_m$ в метке $d$ ненулевой и $\omega(d,x)\ne0\bmod N$, то положителен и переход $v\to x+d$: коэффициент производной содержит множитель $-2\sin(2\pi\bar2\omega(d,x)/N)$. Это доказательство сохраняет точный параметр $t=10^{-80}N$ при любом $m$.

Для $N\ge165$ коэффициенты $D_m$ в $d=(1,0),(0,1)$ ненулевые для всех $m$. Орбита $S^{-j}d$ сохраняет $a^2-2b^2\bmod N$. На исходном носителе целые нормы лежат в $[-162,1]$. Равенство по модулю здесь означает целое равенство, а каждая из норм $1,-2$ встречается только у пары $d,-d$. У первой пары вещественные части сырых коэффициентов равны $-12/4000$, у второй $200/4000$. Член $j=0$ присутствует и отмена невозможна. Сопряжённые метки $-d$ также присутствуют.

Для $3\le N<165$, $N\ne7$, доказательство дополнено конечным целочисленным сертификатом. Используются те же базисные метки, кроме $N=49$, где берутся $(1,1),(0,1)$. Для каждой метки перечисляем одну обратную орбиту длины $h$, складывая исходные гауссовы целые с одинаковыми остатками. Пусть $s_r$ --- сумма первых $r$ коэффициентов, $1\le r\le h$. Для всех $m=kh+r$ нужно проверить $s_r+ks_h\ne0$ при целом $k\ge0$. При $s_h=0$ проверяется $s_r\ne0$; иначе одна ненулевая координата $s_h$ определяет единственное рациональное $k$, после чего проверяются его знак, целость и вторая координата. Всего проверено 7722 префикса. Скрипт \texttt{orbit\_graph.py} воспроизводит весь сертификат из 32 сырых коэффициентов.
\newpage
\sect{14. Строгое сжатие и арифметика концентрации}
Дополним доказательство (16). Обозначим две выбранные метки через $u_1,u_2$; их определитель равен единице. В графе симметричной матрицы $Q Q^{\mathsf T}$ соединены $x$ и $x\pm u_i$ всякий раз, когда $\omega(u_i,x)\ne0\bmod N$, так как обе строки имеют положительный элемент в колонке $P^{-1}x$. В координатах по базису $u_1,u_2$ можно произвольно менять первую координату при ненулевой второй и вторую при ненулевой первой. Из каждой ненулевой вершины можно прийти в $(1,1)$, в том числе при составном $N$. Диагональ $Q Q^{\mathsf T}$ положительна. Связность и дважды стохастичность дают простоту собственного числа 1 по формуле
\[
 \langle f,(I-QQ^{\mathsf T})f\rangle
 =\tfrac12\sum_{x,y}(QQ^{\mathsf T})_{xy}|f(x)-f(y)|^2 .
\]
Сама матрица $QQ^{\mathsf T}$ положительно полуопределена. Тем самым доказано (16). Для любого собственного непустого подмножества $J$ вектор $\mathbf1_J-|J|/(N^2-1)$ имеет строго меньшую норму после действия $Q$; это исключает и точный перенос одного координатного подпространства на другое той же размерности. Все степени $Q^k$ при фиксированном $Q$ сходятся к равномерному усреднению.

Имеется усиление для вероятностей между группами мод. Пусть $A_0,B_0$ --- произвольные подмножества ненулевых меток. Число $H(t)=\sum_{w\in A_0,v\in B_0}|(\operatorname{Ad} U(t)^m)_{wv}|^2$ является экспоненциальным полиномом; $H(0)=|A_0\cap P B_0|$ целое. Если $H(10^{-80}N)$ алгебраично, то $H$ тождественно равно этому целому. В частности, для каждой моды $v\ne0$ вероятность перехода в $Pv$ строго между нулём и единицей: у $Pv$ есть хотя бы один разрешённый сосед. Эта вероятность трансцендентна. Всякая другая положительная вероятность в той же колонке также трансцендентна, поскольку при нуле она обращалась в нуль. Так доказано, что \emph{все} положительные элементы $Q$ трансцендентны.

Разрыв $1-\sigma_2(Q)$ в (16) положителен для каждой степени. При фиксированном $N$ инфимум по $m$ равен нулю. Действительно, по одновременной аппроксимации собственных фаз некоторыми степенями $U_N$ можно сколь угодно близко вернуться к скаляру, и тогда $Q$ приближается к единичной матрице. В частности, аргумент отделяет точную арифметическую изоляцию мод от вопроса о почти изоляции. При переходе от $m$ к $2m$ интерференция остаётся в произведении амплитуд; близость $Q^k$ к усреднению не даёт формулы для $| (W_N^{km})_{wv}|^2$.
\newpage
\sect{15. Все алгебраические резонансы собственных фаз}
Пусть $o_N$ --- порядок $S$ в $\mathrm{SL}(2,\Z/N\Z)$, $\lambda_1,\ldots,\lambda_N$ --- собственные числа $U_N$. Для любых целых $e_j$ с $\sum_j e_j=0$ имеет место
\[
 z=\prod_j\lambda_j^{e_j}\in\overline{\mathbb Q}
                 \quad\Longrightarrow\quad z^{o_N}=1 .                  \tag{17}
\]
В частности, если две собственные фазы различаются на рациональную долю оборота, её приведённый знаменатель делит $o_N$. Для \emph{любого} $m\ge1$ разбиение помеченных собственных значений по равенству $\lambda_i^m=\lambda_j^m$ совпадает с разбиением при степени $\gcd(m,o_N)$. При $\gcd(m,o_N)=1$ степень сохраняет точно все исходные кратности.

Докажем более точную одновременную специализацию. Можно занумеровать собственные числа $\mu_1,\ldots,\mu_N$ линейной карты $M_N$ так, чтобы каждое полиномиальное соотношение над $\overline{\mathbb Q}$, выполненное для всего кортежа $\lambda$, было выполнено также для кортежа $\mu$ в \emph{одной и той же} нумерации. Допустим обратное. Для каждой из конечного числа перестановок $\pi$ найдётся многочлен $P_\pi$, нулевой в $\lambda$, но ненулевой в $\pi\mu$. Можно взять одну линейную комбинацию $P$ этих многочленов с целыми коэффициентами, ненулевую во всех $\pi\mu$: конечное число собственных гиперплоскостей не покрывает целую решётку (либо подставим степени одного целого числа). Тогда $\prod_\pi P(\pi X)$ --- симметрический многочлен. Подстановка собственных значений $U(t)$ выражает его многочленом от коэффициентов характеристического многочлена, то есть конечным экспоненциальным полиномом указанного выше типа. Он равен нулю при заданном алгебраическом $t$ и потому равен нулю при $t=0$. Это противоречит выбору $P$ и доказывает специализацию. Довод работает с кратными собственными значениями, без выбора аналитических ветвей.

Очищаем в $\prod\lambda_j^{e_j}=z$ отрицательные степени. Специализация даёт то же равенство для $\mu_j$. По ковариантности $M_N^{o_N}$ коммутирует со всеми $T(v)$ и потому скалярно. Все отношения $\mu_i/\mu_j$ лежат в группе корней степени $o_N$; ввиду $\sum e_j=0$ получаем (17). Более того, ту же аргументацию можно провести между любыми двумя ненулевыми вещественными алгебраическими значениями $t_1,t_2$. Выбор нумерации в одну сторону и в другую даёт равенство идеалов соотношений с точностью до перестановки: включения по кругу дают включение идеала в его перестановку, которое превращается в равенство после конечного числа итераций.

Для всякого $N\ne7$ среди отношений собственных значений есть хотя бы одно трансцендентное. В противном случае (17) сделало бы $U_N^{o_N}$ скаляром вопреки (16). Тем самым группа отношений зафиксированной матрицы содержит элемент бесконечного порядка, и вся её часть, состоящая из алгебраических чисел, вкладывается в конечную группу $\mu_{o_N}$. Теорема контролирует точные равенства между фазами всех порядков. Для слабого предела $\mathrm{Sine}_2$ существенна частота \emph{приближённых} совпадений в окнах длины $1/N$; (16) и (17) этой частоты не оценивают.

\newpage
\sect{16. Запрещённые резонансы в долях средней щели}
Из (17) следует утверждение на самом масштабе развёртки исходной задачи. Пусть $\lambda_j=e^{2\pi i\theta_j}$, $a_j\in\Z$, $\sum_j a_j=0$ и $s=N\|\sum_j a_j\theta_j\|_{\R/\Z}$. Для всех нечётных $N$ одновременно
\[
 0<s<\tfrac34\quad\Longrightarrow\quad
                   \prod_j\lambda_j^{a_j}\ \text{трансцендентно}.         \tag{18}
\]
Это касается как одной щели, так и разности двух щелей. В границу не входит $s=0$, где возможны точные равенства. В действительности нижний порог для ненулевых алгебраических резонансов равен \emph{по меньшей мере} $N/o_N$; реализация граничного значения в возмущённом спектре не предполагается.

Для доказательства факторизуем $N=\prod_{p\mid N}p^{k_p}$, $r=\#\{p:p\mid N\}$, $\chi_p=(2/p)$. На $a+b\sqrt2$ матрица $S$ действует умножением на $3+2\sqrt2$ нормы 1. Над $\mathbb F_p$ при $\chi_p=1$ две координаты взаимно обратны и порядок делит $p-1$. При $\chi_p=-1$ элемент лежит в ядре нормы в $\mathbb F_{p^2}^*$ и порядок делит $p+1$. Если $S^h\equiv I\pmod p$, то биномиальная формула последовательно поднимает это до $S^{hp^{k-1}}\equiv I\pmod {p^k}$. Китайская теорема об остатках и чётность всех чисел $p^{k_p-1}(p-\chi_p)$ дают
\[
 o_N\le\mathrm{lcm}_{p\mid N}\big[p^{k_p-1}(p-\chi_p)\big]
 \le 2^{1-r}N\prod_{p\mid N}(1-\chi_p/p)
 \le 2(2/3)^r N\le\tfrac43 N .
\]
Таким образом алгебраическое произведение, отличное от 1, имеет $s\ge N/o_N\ge3/4$. При $r\ge2$ уже $s\ge9/8$. Вдоль любой последовательности с $r\to\infty$ порог $N/o_N$ уходит в бесконечность.

Есть и счётное следствие для \emph{соседних} собственных фаз в циклическом порядке, с учётом кратностей. Положительная ориентированная дуга с алгебраическим отношением концов имеет длину хотя бы $1/o_N$ (угловая окружность имеет длину 1). У $N\ne7$ есть хотя бы одна дуга с трансцендентным отношением: иначе все отношения телескопированием были бы алгебраичны, вопреки разделу 15. Сумма длин остальных дуг строго меньше 1. Следовательно, среди $N$ соседних щелей число строго положительных щелей с алгебраическим отношением концов не больше $o_N-1$. При $r\to\infty$ их доля стремится к нулю. Возможные нулевые щели этим подсчётом не ограничены. Оценка также не предписывает распределения длин щелей с трансцендентным отношением.
\newpage
\sect{17. Кратность произвольной точной разности фаз}
Для неподвижного $N$ положим $m_U(z)=\#\{(i,j):\lambda_i/\lambda_j=z\}$, включая $i=j$. Совместная специализация из раздела 15 сохраняет \emph{равенство} любых двух отношений даже при трансцендентном общем значении. На конечном множестве отношений тем самым определено отображение $z\mapsto z_0=\mu_i/\mu_j$. Суммирование $m_U(z)$ по прообразу $z_0$ даёт ровно $m_M(z_0)$. Это одновременно даёт для любого $q\ge1$
\[
 \sum_z m_U(z)^q\le\sum_{z_0}m_M(z_0)^q,\qquad
 \max_z m_U(z)\le m_M(1).
\]
Последняя граница использует неравенство Коши--Шварца для свёртки последовательности кратностей с её сдвигом. То же соображение для равенств любых двух произведений проверяет соответствующее разбиение кортежей всех фиксированных порядков.

Число $m_M(1)$ в точности равно количеству циклов $S$ на $(\Z/N\Z)^2$, поскольку $\operatorname{Ad}M$ переставляет вейлев базис. Пусть $\Lambda=3+2\sqrt2$. Для каждого $h\ge1$ число точек периода, делящего $h$, не превосходит $|\det(S^h-I)|=\Lambda^h+\Lambda^{-h}-2<\Lambda^h$. Действительно, в целой нормальной форме Смита с ненулевыми диагоналями $d_1,d_2$ размер ядра по модулю $N$ равен $\gcd(N,|d_1|)\gcd(N,|d_2|)\le|d_1d_2|$. Короткие циклы длины меньше $H$ дают не более $C\Lambda^H$ циклов, длинные --- не более $N^2/H$. Берём $H=\lfloor (2\log N-3\log\log N)/\log\Lambda\rfloor$ для достаточно больших $N$. Итак, равномерно по \emph{всем} $z\in\mathbb T$ и всем нечётным $N\to\infty$
\[
 \frac{m_U(z)}{N^2}
 \le\frac{\log(3+2\sqrt2)/2+o(1)}{\log N}.                              \tag{19}
\]
В частности, (19) контролирует и заранее неизвестное значение самой часто встречающейся разности фаз.

\subt{Что остаётся открытым} Результаты (16)--(19) относятся непосредственно к фиксированной силе $10^{-80}$, всем степеням в (16) и точным резонансам самого спектра в (17)--(19). Никакая их комбинация пока не определяет вероятность наличия двух разных фаз на расстоянии, например, между $0.2/N$ и $0.3/N$. Для полного парного теста (P) всё ещё нужна суммарная интерференция при $m\asymp N$, а слабый предел точечного процесса требует контроля локальных конфигураций произвольного порядка.

\end{document}
